% !TeX program = xelatex
\documentclass[12pt,a4paper]{article}
\usepackage[margin=2cm]{geometry}
\usepackage{fontspec}
\usepackage{amsmath,amssymb,enumitem}
\setmainfont{Times New Roman}
\setlength{\parindent}{0pt}\setlength{\parskip}{6pt}
\begin{document}
\begin{center}{\Large\bfseries Đề thi Toán vào lớp 10 Hà Nội năm 2014\\Phần Đại số}\end{center}
\textbf{Câu 26.} Tính giá trị của biểu thức
\[
A=\frac{\sqrt{x}+1}{\sqrt{x}-1}
\]
khi $x=9$.

\textbf{Câu 27.} Cho biểu thức
\[
P=\left(\frac{x-2}{x+2\sqrt{x}}+\frac{1}{\sqrt{x}+2}\right):\frac{\sqrt{x}-1}{\sqrt{x}+1},
\qquad x>0,\ x\ne1.
\]
\begin{enumerate}[label=\alph*)]
\item Chứng minh rằng $P=\dfrac{\sqrt{x}+1}{\sqrt{x}}$.
\item Tìm $x$ để $2P=2\sqrt{x}+5$.
\end{enumerate}

\textbf{Câu 28.} Một phân xưởng theo kế hoạch phải sản xuất $1100$ sản phẩm trong một số ngày quy định. Do mỗi ngày phân xưởng sản xuất vượt mức $5$ sản phẩm nên phân xưởng hoàn thành kế hoạch sớm hơn thời gian quy định $2$ ngày. Hỏi theo kế hoạch, mỗi ngày phân xưởng phải sản xuất bao nhiêu sản phẩm?

\textbf{Câu 29.} Giải hệ phương trình
\[
\begin{cases}
\dfrac{4}{x+y}+\dfrac{1}{y-1}=5,\\[3pt]
\dfrac{1}{x+y}-\dfrac{2}{y-1}=-1.
\end{cases}
\]

\textbf{Câu 30.} Trên mặt phẳng tọa độ, cho đường thẳng $(d):y=-x+6$ và parabol $(P):y=x^2$.
\begin{enumerate}[label=\alph*)]
\item Tìm tọa độ giao điểm của $(d)$ và $(P)$.
\item Gọi $A,B$ là hai giao điểm đó. Tính diện tích tam giác $AOB$, với $O$ là gốc tọa độ.
\end{enumerate}

\textbf{Câu 32.} Với $a,b,c>0$ và $a+b+c=2$, tìm giá trị lớn nhất của
\[
Q=\sqrt{2a+bc}+\sqrt{2b+ca}+\sqrt{2c+ab}.
\]
\end{document}
